
\documentclass{article}
\usepackage{colortbl}
\usepackage{makecell}
\usepackage{multirow}
\usepackage{supertabular}

\begin{document}

\newcounter{utterance}

\centering \large Interaction Transcript for game `cladder', experiment `full\_v1.5\_default', episode 9560 with Qwen3.5{-}9B.
\vspace{24pt}

{ \footnotesize  \setcounter{utterance}{1}
\setlength{\tabcolsep}{0pt}
\begin{supertabular}{c@{$\;$}|p{.15\linewidth}@{}p{.15\linewidth}p{.15\linewidth}p{.15\linewidth}p{.15\linewidth}p{.15\linewidth}}
    \# & \multicolumn{2}{c}{Player} && \multicolumn{2}{c}{Game Master} \\
    \hline

    \theutterance \stepcounter{utterance}  
    & & & \multicolumn{4}{p{0.6\linewidth}}{
        \cellcolor[rgb]{0.9,0.9,0.9}{
            \makecell[{{p{\linewidth}}}]{
                \texttt{\tiny{[P1$\langle$GM]}}
                \texttt{You are an expert in causal inference. The following question is not a typical commonsense query, but rather a meticulously designed question created by a professor specializing in causal inference, intended to assess the students' mastery of the course content.} \\
\texttt{Imagine a self{-}contained, hypothetical world with only the following conditions, and without any unmentioned factors or causal relationships: Tijv has a direct effect on xevo and gyzp. Xevo has a direct effect on gyzp. We know that tijv causes not xevo. tijv and xevo causes gyzp. We observed an individual is not tijv. Would an individual is gyzp if not xevo instead of xevo?} \\
\texttt{Start your answer with 'Yes' or 'No', followed by additional reasoning or evidence to support your explanation.} \\
            }
        }
    }
     \\ \\

    \theutterance \stepcounter{utterance}  
    & \multicolumn{4}{p{0.6\linewidth}}{
        \cellcolor[rgb]{0.561,0.737,0.561}{
            \makecell[{{p{\linewidth}}}]{
                \texttt{\tiny{[P1$\rangle$GM]}}
                \texttt{No} \\
\\ 
\texttt{**Reasoning:**} \\
\\ 
\texttt{Based on the conditions provided in this hypothetical world, we can reconstruct the causal structure as follows:} \\
\texttt{1.  **Tijv** has a direct effect on **xevo**.} \\
\texttt{2.  **Tijv** has a direct effect on **gyzp**.} \\
\texttt{3.  **Xevo** has a direct effect on **gyzp**.} \\
\\ 
\texttt{The prompt states: "We know that tijv causes not xevo." This establishes a deterministic or probabilistic constraint where the state of `Tijv` determines the state of `Xevo`. Specifically, if `Tijv` is present (True), then `Xevo` is absent (False). Conversely, if `Tijv` is absent (False), the prompt does not explicitly state the value of `Xevo`, but it implies `Xevo` is not *caused* to be false by `Tijv`. However, the critical part of the query is the counterfactual: "Would an individual be gyzp if not xevo instead of xevo?" given that the individual is "not tijv".} \\
\\ 
\texttt{Let's analyze the counterfactual scenario:} \\
\texttt{*   **Factual World**: The individual is **not Tijv**.} \\
\texttt{    *   Since `Tijv` is false, the direct causal path from `Tijv` to `gyzp` is inactive (or contributes nothing).} \\
\texttt{    *   The state of `Xevo` is unknown from the text alone, but the question asks us to evaluate the outcome if we intervene to make `Xevo` false ("if not xevo").} \\
\texttt{*   **Counterfactual Intervention**: We set **Xevo = False** (not xevo).} \\
\texttt{*   **Causal Path Analysis**:} \\
\texttt{    *   The variable `gyzp` has two parents: `Tijv` and `Xevo`.} \\
\texttt{    *   In this specific instance, `Tijv` is False.} \\
\texttt{    *   The intervention sets `Xevo` to False.} \\
\texttt{    *   Therefore, both direct causes of `gyzp` are False.} \\
\\ 
\texttt{Unless there is an unmentioned background cause or a default state (which the prompt explicitly forbids by stating "without any unmentioned factors"), `gyzp` will not be caused to occur if both its direct causes (`Tijv` and `Xevo`) are absent. The prompt states `Tijv` causes `gyzp` and `Xevo` causes `gyzp`. If neither is active, `gyzp` does not occur.} \\
\\ 
\texttt{Therefore, if the individual is not `Tijv` and we intervene to make them not `Xevo`, the result is that `gyzp` does not happen. The answer to "Would an individual be gyzp?" is No.} \\
            }
        }
    }
    & & \\ \\

    \theutterance \stepcounter{utterance}  
    & & & \multicolumn{2}{p{0.3\linewidth}}{
        \cellcolor[rgb]{0.9,0.9,0.9}{
            \makecell[{{p{\linewidth}}}]{
                \texttt{\tiny{[GM$|$GM]}}
                \texttt{No} \\
            }
        }
    }
    & & \\ \\

    \theutterance \stepcounter{utterance}  
    & & & \multicolumn{2}{p{0.3\linewidth}}{
        \cellcolor[rgb]{0.9,0.9,0.9}{
            \makecell[{{p{\linewidth}}}]{
                \texttt{\tiny{[GM$|$GM]}}
                \texttt{no} \\
            }
        }
    }
    & & \\ \\

    \theutterance \stepcounter{utterance}  
    & & & \multicolumn{2}{p{0.3\linewidth}}{
        \cellcolor[rgb]{0.9,0.9,0.9}{
            \makecell[{{p{\linewidth}}}]{
                \texttt{\tiny{[GM$|$GM]}}
                \texttt{game\_result = WIN} \\
            }
        }
    }
    & & \\ \\

\end{supertabular}
}

\end{document}
